\medskip\noindent\textbf{Supplementary Note 7. Secondary analyses.}\par\smallskip
\begingroup\small\setlength{\parskip}{3pt}
These descriptive analyses use seed-mean batch and biology scores $B,C$, with plain-scVI means $b_0,c_0$. Supplementary Table S8 reports results by penalty, task and decoder.

\noindent\textit{Hypervolume.} Each penalty strength defines a rectangle from $(b_0,c_0-0.10)$ to $(B,C)$ when both scores exceed that reference. Hypervolume is their union's area divided by $b_{\mathrm{common}}-b_0$. Batch scores are capped at $b_{\mathrm{common}}$ before forming rectangles. Plain scVI contributes no rectangle. No qualifying point gives 0, and low-signal pairs are excluded. On the scaled batch axis $x=(B-b_0)/(b_{\mathrm{common}}-b_0)$, this is area beneath the Pareto frontier over $x=0$ to 1, above biology difference $-0.10$. Retaining plain scVI's biology throughout gives 0.10.

The uncapped value, HV$_{\mathrm{all}}$, also counts batch removal beyond the common range. The barycenter critic's maximum set $b_{\mathrm{common}}$ in all 11 ranked real pairs, while every other penalty reached further. The uncapped measure therefore rewards their additional range. Figure 3 shows frontiers and Supplementary Figure S8 shows capped hypervolumes.

\noindent\textit{Non-dominated strengths.} A point is dominated if another strength of the same penalty is at least as good on both scores within $10^{-9}$, and better on one by more than $10^{-9}$. Identical points do not dominate each other. The share divides non-dominated points by all points, omitting strengths without successful fits. Low-signal pairs remain eligible.

\noindent\textit{Collapse.} A fit collapses if it fails or its biology falls at least 0.10 below seed-matched plain scVI, within $10^{-9}$. The main-text rate uses strengths whose mean batch score is at most $b_{\mathrm{common}}$, including strong settings that fall back below it. Supplementary Table S8 also reports all-strength rates. All pairs are included. No sweep fit failed, so every collapse reflects biology loss. Exact binomial (Clopper--Pearson) 95\% intervals are nominal because fits sharing a seed also share their reference threshold.

An illustrative simulation used seed 7 and 20,000 replicates per scenario for grids of ten or twelve strengths with three seeds. Independent collapse probabilities followed the profile 0.0, 0.0, 0.05, 0.2, 0.6, 0.95, linearly interpolated to each grid. A dependent scenario instead assigned one threshold per seed, after which all higher strengths collapsed. Threshold probabilities were 0.0, 0.0, 0.0, 0.1, 0.15, 0.2, 0.55 across seven positions, the last meaning no collapse, interpolated and renormalised. An independent control used probability 0.3 and required coverage at least 0.95. Coverage is the share of replicates whose 95\% interval contained the grid's mean collapse probability. It was at least 0.999 with independent varying probabilities, 0.88 (30 fits) and 0.83 (36 fits) with shared thresholds, and 0.97 and 0.95 in the control. Monte Carlo standard errors were at most 0.003. These chosen profiles demonstrate possible undercoverage, without bounding coverage in the sweep. The 33-fit design was not simulated.

\noindent\textit{Single-strength placement.} For each real task, the transferred strength maximises mean selection score $S=0.4B+0.6C$ on real tasks of other families. The four families are pancreas; lung; the immune and human--mouse immune tasks; and the two chromatin-accessibility tasks. Eligible strengths must be fitted and failure-free in every training task and fitted in the target task. Ties favour smaller strengths. Simulations use strengths chosen on all real tasks. The task's own best failure-free strength is an optimistic bound. Harmony, Scanorama and sysVI use the same selection rule (Supplementary Note 8). Supplementary Tables S7 and~\ref{tab:4} compare these scores with plain scVI.
\par\endgroup
